% =====================================================================
\chapter{反对意见与限度}\label{ch:objections}
% =====================================================================

本章执行一项自查：把不可承重的主张移出结构，并显式记录移出了什么。

\section{不作主张清单}

以下清单是硬性的：本文的任何其他部分都不得被解读为支持下列主张。

\begin{enumerate}[label=(N\arabic*)]
\item 本文\emph{不}主张所有大型软件最终收敛到同一种高层架构（假设
      \ref{hyp:strong} 被明确拒绝）。
\item 本文\emph{不}主张代码量本身导致内部基础设施化；本文主张的是实体数与资源数的
      乘积、以及重复问题与演化频率。
\item 本文\emph{不}主张六职能必须\emph{自建}；本文主张的是它们必须被\emph{定义}，
      可以外借给宿主。
\item 本文\emph{不}主张 Simon 的近可分解性是一条普遍定律。本文把它改造为一个有条件、
      可检验的充分统计量判据（定义 \ref{def:neardecomp}），并承认原文献不提供验证
      成本的普遍定理。
\item 本文\emph{不}主张最小描述长度可以决定系统应当实现什么目的。MDL 是模型选择
      原则，不是意图理论。
\item 本文\emph{不}主张"理解就是压缩"在所有领域都是定理。它在归纳推断范围内是可
      辩护的理论立场，在超出该范围时是有力的启发式。
\item 本文\emph{不}主张机械系统绝不犯错。本文主张的是：机械判定不依赖生成者的
      主观意愿；机械系统仍然可能有缺陷、故障与不完整的模型。
\item 本文\emph{不}主张信任可以被完全机械化。TCB 可以被缩小、验证与隔离，但不能被
      凭空消除。
\item 本文\emph{不}主张意图层可以被机械封闭。它可以被对抗性共识\emph{逼近}，
      且逼近的可靠性受相关错误模式支配。
\item 本文\emph{不}主张"加更多智能体来审查"能提高可靠性。按推论
      \ref{cor:diversity}，只有降低相关错误模式才能提高可靠性。
\item 本文\emph{不}主张分层可以容纳一切。横切关注点在结构上不能被局部化到任何一层
      （原则 \ref{pr:crosscutting}）。
\item 本文\emph{不}主张形式化验证会在 AI 时代\emph{自动}变得便宜。这是一个
      值得研究的假说，而不是已成立的事实（第 \ref{sec:asymmetry} 节）。
\item 本文\emph{不}主张编译正确性等于系统正确性。链接、装载、运行时、外部库、
      硬件与操作系统未必在证明范围内。
\item 本文\emph{不}主张资源保证可以统一为一种方法。时间复杂度、符号资源上界、
      最坏执行时间、统计分位数与实际测量是不同等级的证据。
\item 本文\emph{不}主张"未检测到违反"等于"所有执行均满足"。三个状态必须被区分
      （原则 \ref{pr:three-states}）。
\end{enumerate}

\section{被移出承重结构的论证}

本文主动拒绝若干在同类讨论中常见、但不成立的论证。拒绝它们不是因为它们无用，
而是因为它们把启发式写成了定理，从而使整个框架依赖未经证明的猜想。

\begin{longtable}{L{3.4cm}L{5.2cm}L{5.6cm}}
\caption{被移出承重结构的论证}\label{tab:overreach}\\
\toprule
\textbf{被拒绝的论证} & \textbf{为什么不成立} & \textbf{本文的替代} \\
\midrule
\endfirsthead
\toprule
\textbf{被拒绝的论证} & \textbf{为什么不成立} & \textbf{本文的替代} \\
\midrule
\endhead
\bottomrule
\endlastfoot
由 Gödel 不完备定理推出"意图与规格之间必然存在永久裂缝" &
  不完备定理讨论特定形式系统中的可证明性限制，不是关于人的心理意图与形式规格之间
  关系的定理 &
  用三条认识论根据替代：需求是生成性的；规格是有损压缩；验证与确认的区分
  （命题 \ref{prop:three-grounds}） \\
由 Rice 定理推出"意图无法形式化" &
  Rice 定理讨论程序的非平凡语义性质的不可判定性，与人的意图无关 &
  只把它用作通用综合与通用判定算法的边界（注 \ref{rmk:q1q8} 的 Q6） \\
把"判卷比答卷容易"称为 $\mathrm{P}\neq\mathrm{NP}$ 的红利 &
  $\mathrm{P}$ 是否等于 $\mathrm{NP}$ 仍开放；该论证依赖未证明猜想 &
  只使用 $\mathrm{NP}$ 类定义所蕴含的结构性不对称（命题
  \ref{prop:asym-struct}） \\
断言"验证器必须极小、确定、绝不同构"是数学定理 &
  小型可信内核是有价值的架构原则，小不等于绝无缺陷 &
  TCB 最小化：让不可信部分可被检查，让不可避免的信任集中在尽可能小、性质尽可能
  明确的基础上 \\
断言"死的东西不会背叛"，故机械保证绝对可靠 &
  机械系统可能有缺陷、故障与不完整的模型 &
  区分"不背叛"（不依赖生成者意愿）与"无误"（不可能） \\
断言 Simon 层级定律证明所有复杂系统必然层级化 &
  原文献提供解释，不提供普遍定律，也不给出验证成本的普遍定理 &
  改造为有条件、可检验的充分统计量判据（定义 \ref{def:neardecomp}） \\
断言线性逻辑使资源浪费在结构上不可能 &
  线性逻辑的指数结构仍允许特定形式的复制与忽略；线性使用不等于最优使用 &
  采用组合结构：使用约束 $+$ 资源上界分析 $+$ 性能模型 $+$ 运行时计量 $+$
  符合条件的最坏执行时间分析 \\
断言形式化验证在 AI 时代必然变得便宜 &
  历史成本量级显示规格设计、模型选择、证明架构仍是实际工作（例
  \ref{ex:sel4-cost}） &
  降为假说，并给出可测量的成本指标（第 \ref{ch:program} 章判据 S5） \\
断言"理解即压缩"是普遍定律 &
  其严格形式化适用于归纳推断，超出该范围退化为启发式 &
  限定适用范围，并把它与任务相关的充分性绑定（定义 \ref{def:sufficient}） \\
断言本框架与 Langlands 纲领存在实质联系 &
  该纲领研究算术对象与自守表示之间的函子性对偶；本框架所需的"生成自由 / 准入确定"
  只是\emph{对偶的气味}，不是函子。写入它会以修辞替代结构 &
  采用能付账的对偶：LP/SMT 对偶。Farkas 证书、不可满足核与不可实现证书给出了
  综合问题与其不可行性证明之间的定理级对应，且已在工程中使用 \\
断言 Landauer 原理约束软件设计 &
  软件尺度与该下界相差约十二个数量级；擦除下界不是架构的约束面 &
  采用存储墙与局部性：每比特搬运的能耗比擦除下界高约 $10^{9}$ 倍，这才是决定内存
  层级、缓存策略与数据布局的实际物理 \\
断言同伦类型论是本框架的必要基础 &
  该理论扩展了数学基础，但未给出资源分析或已验证编译尚未提供的软件保证 &
  将类型论的作用限于已被证明的部分：依赖类型表达后置条件与不变量
  \cite{hott2013} \\
\bottomrule
\end{longtable}

\section{可能的反对意见与回应}

\paragraph{反对意见一：读集与无视集在真实系统中无法精确测量，因此这些形式化是
装饰性的。}
回应：定义 \ref{def:readset} 的读集可以从版本历史中以"每次变更触及的模块集合"作为
代理直接统计，这是一个可计算的量（假设 $H_4$ 正是要检验它作为代理的有效性）。
无视集在本文中并非要求被自动提取，而是要求被\emph{显式声明}，它是设计产物而非分析
产物。声明之后，它可以被泄漏对诊断程序（推论 \ref{cor:diagnose}）部分地检验。

\paragraph{反对意见二：核心结论"大系统会长出一个内核"过于宽泛，
几乎无法证伪。}
回应：定理 \ref{th:six} 有明确前提（并发实体多于一个、资源有限且受争用、故障可达），
且其结论是可证伪的：给出一个满足前提却缺少六项之一的良定义系统即可推翻它。
经验部分则由假设 $H_3$ 承担，其证伪条件是"代码行数的解释力不劣于 $n \times m$"。

\paragraph{反对意见三：双层信任结构把太多责任推给共识，而共识在实践中常常失效。}
回应：这正是命题 \ref{th:correlated} 要处理的问题。该命题不是为共识辩护，而是给出
它的失效形式与设计变量。它同时解释了一个常见现象：由同源模型组成的"审查团"在结构上
不构成保证，无论其规模多大。

\paragraph{反对意见四：把层级当作验证义务的分解单元，是否只是在重新命名已有的
"分层验证"实践？}
回应：部分是的，本文承认这一点。区别在于三点：(a) 本文给出的归属规则
（定义 \ref{def:obligation-assign}）是一个可操作的判定，而不是一个口号；
(b) 本文显式引入了"该层不可表达的义务"这一类，并要求它被分派给共识或人的审定，
而不是被忽略；(c) 本文给出共识层可靠性的定量界，从而使"分派给共识"成为一个可以被
评估而非只能被信任的决定。

\paragraph{反对意见五：把技术债定义为描述长度差是不可计算的，因此没有用。}
回应：该定义的价值不在于数值计算，而在于它给出的方向与判据（原则
\ref{pr:debt}(c)）：判断一次重构是否值得，等价于判断它是否让系统更容易被更短的描述
生成。这是一个可以在评审中被实际使用的问题，即使 $\DL^{\star}$ 不可计算。

\paragraph{反对意见六：本文混合了工程、数学与社会认识论，是否超出了单一学科的可
辩护范围？}
回应：这一批评成立。缓解措施有两项。第一，对每一类主张标注其强度
（第 \ref{ch:program} 章的三分法）；第二，明确区分哪些引用是被\emph{使用}的、
哪些是被\emph{拒绝}的（表 \ref{tab:relations}）。跨学科本身不是缺陷，前提是每个
学科的结论都不被越界使用。

\section{尚未解决的问题}

以下问题是开放的：

\begin{enumerate}[label=(O\arabic*)]
\item \term{$\mathrm{Dom}$ 与 $\mathrm{Conv}$ 的独立可测性}。猜想
      \ref{conj:strength} 要成为可检验命题，需要两个量都能被独立于结论地度量。
      本文只给出了方向，没有给出可复现的标注规程。
\item \term{六职能完备性论证的形式化}。定理 \ref{th:six} 的证明是结构性枚举，
      不是形式系统中的推导。把它写成对一个给定的操作语义框架的完备性定理，是一项
      尚未完成的工作。
\item \term{跨层语义的组合是否总能成立}。命题 \ref{prop:soundness} 给出充分条件，
      但在包含非确定性、实时约束与自反性的系统中，这些条件是否可同时满足尚不清楚。
\item \term{共识层的 $q$ 如何度量}。命题 \ref{th:correlated} 把可靠性归因于相关错误
      模式，但没有给出 $q$ 的可操作估计方法。这是该方向能否落地的关键。
\item \term{资源与功能的联合语义是否可判定}。带成本的迹语义可以表达联合约束，
      但其可判定片段有多大、紧致度如何，尚无定论。
\item \term{制品绑定与动态生成实现的兼容}。如果实现可以在运行时根据负载重新生成，
      式 \eqref{eq:binding} 的静态同一性断言需要被替换为"哪些变化保持原有契约、
      哪些要求重新验证、验证失败时如何回退到已知安全状态"这一更强的问题。
\item \term{层级归属的一致性}。定义 \ref{def:obligation-assign} 是否能在不同分析者
      之间给出一致的归属，需要实证检验。
\end{enumerate}

\section{方法论局限}

\begin{enumerate}[label=(L\arabic*)]
\item 本文引用的全部文献，其核验深度限于正式出版页面、作者或项目主页、以及预印本
      元数据，\emph{未逐篇通读全文}。因此文中每一个具体数字都应被视为
      \emph{待复核的二手报告}；标称年份较新的若干条目按原始记录保留，本文未做独立
      核验。
\item 本文的形式化是\emph{为本文目的而建立的}工作定义，不应被误称为已有统一标准。
      特别是"无视集""层契约四元组""验证义务的层级归属"三项，是本文提出的工作
      概念。
\item 本文的三条认识论根据（命题 \ref{prop:three-grounds}）是论证而非定理；
      其可挑战之处在于这三条是否穷尽。
\item 本文未进行任何实验。第 \ref{ch:program} 章的全部假设均处于"待检验"状态。
\end{enumerate}

\begin{principle}[自查结论]\label{pr:self-audit}
上述审视之后，承重结构剩下三条可独立检验的最小命题：
\begin{enumerate}[label=(\roman*)]
\item 验证与生成之间存在由 $\mathrm{NP}$ 类定义所蕴含的结构性不对称；
\item 层级是保证义务的分解单元，因为每条义务都归属某个可观察域；
\item 意图与规格之间的鸿沟有认识论根据，因而不可被普遍机械封闭。
\end{enumerate}
其余主张（包括"压缩""涌现""内部基础设施化"等）要么已被降为可证伪假设，
要么已被限定在明确的适用范围内。
\end{principle}
